In platform engineering, \enquote{capabilities} are often the center of discussion around a platform.
A capability here is a way of describing a service a platform offers.
A capability can be as simple as \enquote{runs a container} and as complex as \enquote{provisions a clustered database across regions, with constant \gls{wal} archiving to achieve \gls{pitr}}.
The promises explained in previous chapter map closely to capabilities.
\parencite{platformengineering-architects, cncfPlatformEngineering}
\par
The platform offers a service to a user, which can be directly mapped to a promise \parencite{promise-theory}:

\[ P \xrightarrow{+b} U \]

Where $P$ is the platform and $U$ is a user of the platform.
This means our capability would be $b$ in this case.

\subsection{Promising impositions in a platform}
Since Ansible is imposition based and Kubernetes is imposition based, to provide Ansible as a Service using Kubernetes, a possibility needs to be engineered to map an imposition into a promise.
The simplest and most common way to achieve that is by making a promise to impose.
\parencite{promise-theory}

First Ansible $A$ needs to impose any configuration $c$ on any Server $S_?
	\in S$.

\[ \pi _1: A \impose{+c}{5}
	S_?
\]

The server $S_?
$ then makes a promise to accept the imposition by allowing Ansible to access it:

\[ \pi _2: S_? \xrightarrow{-c}
	A \]

This is important, because an imposition on a autonomous agent is only as strong as the promise by the target agent to accept it.
\parencite{promise-theory}
\par
Following that, Ansible can promise its own behavior, since it is itself an agent that may make promises about its own behavior:

\[ \sigma _1: A \xrightarrow{+\pi _1} P \]

This means Ansible makes a promise to the platform $P$ to impose $c$ on any Server $S_?
$ in the list of servers $S$.
The platform promises to use the promise, which allows the platform to depend on a promise made by Ansible.

\[ \sigma _2: P \xrightarrow{-\pi _1} A \]

Additionally, Ansible needs to provide feedback on whether it was successful in imposing upon the server, so the platform can assess whether Ansible has fulfilled its promise:

\[ \phi _1: A \xrightarrow{+f} P, \phi _2: P \xrightarrow{-f} A \]

This is now a promise, however, it is simply Ansible promising its own behavior to our platform.
The issue here is that our platform cannot promise Ansible's behavior, since this would be in violation of the principle of autonomy.
\parencite{promise-theory}
\par
Instead, we may promise an outcome and make it subject to the assessment of another promise:

\[ \rho _1: P \xrightarrow{+D|\alpha _P(\sigma _1)} U \]

This means our platform $P$ now promises to provide a desired state $D$ to our User $U$, subject to a successful assessment $\alpha$ by our platform $P$ of the promise $\sigma _1$.
It should be noted that we explicitly rely on $\alpha _P(\sigma _1)$ based on $f$, the assessment of P on $\sigma _1$, rather than on the server being in the desired state, since the platform $P$ does not observe the state of $S_?
$.
It can merely use the output provided by Ansible to assess whether it kept the promise.
Since it is the body of our promise, our concrete capability is now $D$, which means our platform now has the capability to use a promise Ansible has made to it to apply a configuration to an arbitrary server.
\parencite{promise-theory}
